% GDZ PPN236018647,native512–516;HKUST513–517. Native513 blank.
\clearpage\SourceLeaf{512}\NoSourceHead
\pdfbookmark[1]{Abdruck des Tagebuchs (Notizenjournals) mit Erläuterungen}{diary-printed-title}
\vspace*{.20\textheight}
\begin{center}\Large
ABDRUCK DES TAGEBUCHS\\[5pt]
(NOTIZENJOURNALS)\\[5pt]
MIT ERLÄUTERUNGEN.
\end{center}
\vfill\noindent\hfill61\textsuperscript{*}
% Native513 is blank of historical print.
\clearpage\SourceLeaf{514}\NoSourceHead
\begin{center}\bfseries
VORBEMERKUNGEN ZUM ERSTEN ABDRUCK DES TAGEBUCHS
\end{center}
\noindent in der »\textit{Festschrift zur Feier des hundertfünfzigjährigen Bestehens der Königlichen Gesellschaft der Wissenschaften zu Göttingen}«, Beiträge zur Gelehrtengeschichte Göttingens, Berlin 1901; siehe auch Mathematische Annalen \textbf{57}, 1903, S. 1—34\footnote{Diese \textls[120]{Vorbemerkungen} werden hier mit Weglassung einiger nur für den ersten Abdruck in Betracht kommender Stellen und mit einigen durch die gegenwärtigen Verhältnisse bedingten Abänderungen wiedergegeben.}.

\medskip
\indent Die tagebuchartigen Aufzeichnungen von \textsc{Gauss} füllen in der Urschrift 19 kleine Oktavseiten in einem unscheinbaren Heftchen, welches sich seit \textsc{Gauss}' Tode in der Familie weiter vererbt hat und uns durch Vermittlung von Herrn \textsc{Stäckel} im Sommer 1898 seitens des Enkels von \textsc{Gauss}, des Herrn C. \textsc{Gauss} in Hameln, zur Benutzung bei der Weiterführung der Gesamtausgabe von \textsc{Gauss}' Werken zur Verfügung gestellt wurde. Herr C. \textsc{Gauss} hat später — unter Wahrung seines persönlichen Eigentumsrechts — in dankenswerter Weise verfügt, daß besagtes Heftchen dauernd im hiesigen \textsc{Gauss}archiv aufbewahrt werden soll.

\indent Die außerordentliche wissenschaftliche Bedeutung dieses \textit{Tagebuchs} oder \textit{Notizenjournals} (wie es \textsc{Gauss} selbst gelegentlich in einem Brief an \textsc{Olbers} nennt; siehe Nr. 88 des folgenden Abdrucks) ist bereits im zweiten derjenigen Berichte, welche der K. Gesellschaft der Wissenschaften alljährlich über den Stand der Herausgabe von \textsc{Gauss}' Werken erstattet werden, unter Mitteilung einiger charakteristischer Stellen hervorgehoben worden\begingroup\renewcommand{\thefootnote}{**)}\footnote{Nachrichten der K. Gesellschaft der Wissenschaften zu Göttingen 1899, Geschäftliche Mitteilungen, Heft 1. (Abgedruckt Mathematische Annalen \textbf{53}, 1899, S. 45—48.)}\endgroup, — sie tritt nicht minder in dem Bande VIII von \textsc{Gauss}' Werken hervor, wo wir vielfach auf die Angaben des \textit{Tagebuchs} verweisen konnten. Nach dem von der K. Gesellschaft angenommenen allgemeinen Plane für die Weiterführung der Gesamtausgabe soll dasselbe in Band X der Werke ausführlich publiziert und bearbeitet werden. Aber es ist bis dahin noch ein langer Weg, dessen Ende noch nicht mit Sicherheit abzusehen sein dürfte. Ich glaube also auf allgemeine Zustimmung rechnen zu dürfen, wenn ich das \textit{Tagebuch} hier vorab als Beitrag zu der von der K. Gesellschaft der Wissenschaften anläßlich ihres 150jährigen Bestehens geplanten historischen Festschrift in \textls[120]{vorläufiger Form} veröffentliche. Die endgültige Bearbeitung in Band X wird damit nichts an ihrem Werte verlieren, sie wird aber dadurch, daß das Material schon jetzt zur öffentlichen Kenntnis und öffentlichen Diskussion kommt, erleichtert werden.

\indent Zwei Dinge dürfen ja wohl gleich hier vorab hervorgehoben werden, welche dem \textit{Tagebuch} einen unvergleichlichen biographischen Wert verleihen.

\clearpage\SourceLeaf{515}\ArithmeticHead{486}{L}{TAGEBUCH MIT ERLÄUTERUNGEN.}
\indent Das eine ist der unmittelbare, sozusagen persönliche Einblick, den wir gerade für die entscheidenden Jahre 1796—1800 in den wissenschaftlichen Werdegang des jungen \textsc{Gauss} gewinnen\footnote{\textsc{Gauss} ist am 30. April 1777 geboren, war also, als er das \textit{Tagebuch} am 30. März 1796 begann, noch nicht ganz 19 Jahre alt.}. Da ist noch keine Spur der abgeschlossenen Reife des wissenschaftlichen Urteils oder auch der vornehmen Zurückhaltung, wie sie \textsc{Gauss} in späteren Lebensjahren zu eigen waren. Wichtiges und Unwichtiges wechselt ab; neben Entdeckungen von der größten Tragweite finden sich Trivialitäten, wie sie der Anfänger zu überwinden hat; überall aber tritt der persönliche Anteil, den \textsc{Gauss} an seinen Mühen und Erfolgen nimmt, in überquellender Weise zu Tage. Und dabei immer wieder die Eigenart seines mathematischen Genius: induktiv, an der Hand von Zahlenrechnungen, die Resultate zu finden, um hinterher langsam, in härtester Arbeit, die Beweise zu zwingen.

\indent Das andere ist die Verknüpfung der einzelnen wissenschaftlichen Fortschritte, die \textsc{Gauss} gelingen, und die genaue Datierung bestimmter Entdeckungen. Man kann allerdings den Wunsch nicht unterdrücken, \textsc{Gauss} möchte bei seinen Aufzeichnungen mit größerer Konsequenz vorgegangen sein und sich außerdem nicht so vielfach mit bloßen Andeutungen begnügt haben. Manche Frage nach der Entstehung von \textsc{Gauss}' späteren mathematischen Auffassungen und Ideen wird sich auch mit Hilfe des \textit{Tagebuches} niemals beantworten lassen und andererseits wird uns manche Tagebuchnotiz dauernd unverständlich bleiben. Trotzdem ist der Fortschritt, der sich aus einem genauen Vergleich der einzelnen Nummern des \textit{Tagebuches} mit den erhaltenen Stücken des Nachlasses ergeben muß, zweifellos ein sehr bedeutender. Ich habe bei dieser ersten Publikation als meine Hauptaufgabe angesehen, in dieser Hinsicht die Wege nur erst zu ebnen, und bin nur nach \textls[120]{einer} Seite weiter gegangen, indem ich nämlich das Material aus den Jahren 1797—1800 heranzog, welches sich auf die \textit{Theorie der elliptischen Funktionen} bezieht. Das Hauptergebnis meiner betreffenden Studien findet sich in einer Bemerkung zu Nr. 111 des \textit{Tagebuches}: die Fachgenossen müssen entscheiden, wie weit sie dasselbe als gesicherten Besitz anerkennen und dementsprechend die bisher geltende Auffassung abändern wollen.

\indent Hinsichtlich, der Art der im folgenden gegebenen Veröffentlichung und der hinzugefügten Bemerkungen mögen übrigens folgende Angaben vorausgeschickt werden.

\indent Zunächst, was die Notizen des \textit{Tagebuches} selbst betrifft, so habe ich mich bei deren Wiedergabe keineswegs genau an die Äußerlichkeiten des Originals gebunden. Vor allen Dingen habe ich im Interesse der Übersichtlichkeit und der bequemeren Bezugnahme die sämtlichen Sätze durchlaufend numeriert. Die Angaben über Ort und Zeit wurden möglichst gleichförmig gestaltet, und auch dem Text, wo es wünschenswert schien, hin und wieder ein Wort oder eine Silbe (die dann in eckige Klammern eingeschlossen sind) hinzugefügt. Leicht erkennbare sprachliche Unrichtigkeiten wurden verbessert. Die Formeln wurden herausgehoben und in moderner Weise gedruckt. Einer Anzahl Nummern hat \textsc{Gauss} gewisse Marken vorangestellt, um deren Wichtigkeit hervorzuheben, ferner ist eine große Zahl der Notizen, insbesondere derjenigen zahlentheoretischen Inhalts, im Original mit roter Tinte unterstrichen; im Druck wurden sowohl die Marken wie auch die Unterstreichungen weggelassen. Es scheint, als seien die Unterstreichungen erst hinterher angebracht und darauf bezüglich, ob \textsc{Gauss} die einzelne Notiz bei späteren Arbeiten benutzt hat oder nicht.

\indent Dann aber, was die Bemerkungen angeht, die den einzelnen Nummern zugesetzt sind, so haben sie vorwiegend den Zweck, den aktenmäßigen Wert des mitgeteilten Materials zu erhöhen. Hierzu schien vor allen Dingen der Hinweis auf parallellaufende Zeitangaben in den bisherigen Veröffentlichungen von \textsc{Gauss}' Werken oder Briefen erwünscht. Besonderen Dank habe ich Herrn \textsc{Dedekind} für einige erläuternde Bemerkungen zu sagen, die mit seiner Namensunterschrift den in Betracht kommenden Nummern hinzugefügt worden sind. Hierüber hinaus habe ich verschiedentlich auf den handschriftlichen Nachlaß von \textsc{Gauss}, wie
\clearpage\SourceLeaf{516}\ArithmeticHead{487}{R}{VORBEMERKUNGEN ZUM TAGEBUCH.}
\noindent er zur Zeit im hiesigen \textsc{Gauss}archiv aufbewahrt wird, Bezug genommen und insbesondere bei denjenigen Nummern, die sich auf die elliptischen Funktionen in den Jahren 1797—1800 beziehen, solche Stücke des Nachlasses abgedruckt, die nun erst an der Hand des \textit{Tagebuchs} ihre volle Bedeutung gewinnen. Für die Jahre 1796, 1797 ist in dieser Hinsicht ein mit Schreibpapier durchschossenes Exemplar des Lehrbuchs von \textsc{Leiste}: \textit{Die Arithmetik und Algebra}, Wolfenbüttel 1790, (114 pag.) besonders wertvoll, in dem \textsc{Gauss} damals auf dessen freie Seiten eine Reihe der interessantesten Eintragungen gemacht hat (wie er ja überhaupt in die Bücher seiner Bibliothek vielfach Notizen eintrug, gleich als wollte er jedes leere Blatt ausnutzen, das Dauer zu besitzen schien)\footnote{Die Notizen aus \textsc{Leiste} werden in der Folge so zitiert, daß jedesmal die Druckseite angegeben wird, \textls[120]{neben} der sie sich in dem durchschossenen Exemplare befinden.}. Für die Jahre 1798—1800 kommen dann neben losen Zetteln, die sich zufällig erhalten haben, insbesondere die sogenannten \textls[120]{Schedae} in Betracht, d. h. Notizheftchen, welche in ungeregelter Aufeinanderfolge Zahlenrechnungen und Bemerkungen der verschiedensten Art, vielfach auch die Ansätze zu zusammenhängenden Darstellungen enthalten; das Nähere hierüber ist unten bei den einzelnen Nummern mitgeteilt. —

\indent Ich habe noch nach verschiedenen Seiten hin für vielfache Unterstützung, die ich bei meiner Arbeit fand, Dank auszusprechen. Herr \textsc{Brendel} dahier (der jetzige Generalredaktor der \textsc{Gauss}ausgabe) hat mich durch seine große Kenntnis des Nachlasses weitgehend unterstützt. Nicht minder bin ich den Bearbeitern von Band VIII der \textsc{Gauss}schen Werke, den Herren \textsc{Börsch}, \textsc{Fricke} und \textsc{Stäckel}, sowie den Herren \textsc{Fueter} und \textsc{Sommer} dahier für vielfache Bemerkungen und sonstige Hilfe verpflichtet. Der Mitwirkung von Herrn \textsc{Dedekind} gedachte ich schon oben; sie erstreckte sich schließlich auf fast alle Teile des \textit{Tagebuchs} und ist mir besonders wertvoll gewesen.
\par\indent Göttingen, den 3. Juli 1901.
\par\noindent\hfill F. \textsc{Klein}.

\bigskip
\begin{center}\bfseries
VORBEMERKUNG ZU DEM HIER FOLGENDEN ABDRUCK\\
DES TAGEBUCHS.
\end{center}
\indent Die Auffindung des \textit{Tagebuchs} durch \textsc{Stäckel} und seine erste Herausgabe durch \textsc{Klein} bedeuten die Eröffnung eines neuen Abschnitts in der \textsc{Gauss}forschung, die durch diesen Stoff neubelebt und in die Bahnen zuverlässiger geschichtlicher Untersuchung geleitet worden ist. Die nachfolgende vom Unterzeichneten besorgte Ausgabe schließt sich eng an die erste an, nur konnten die Bemerkungen vermehrt und ausführlicher gestaltet werden; namentlich bei solchen Aufzeichnungen, die sonst nirgendwo berührte Gegenstände betreffen, haben sie oft die Form von kleinen Abhandlungen angenommen. Die einzelnen Bemerkungen sind mit den Namen ihrer Urheber unterzeichnet; der bei vielen Nummern diesen Namen vorangestellte Name \textsc{Kleins} soll darauf hinweisen, daß die betreffende Aufzeichnung in der ersten Ausgabe erläutert war, und daß der wesentliche Inhalt jener Erläuterung in die jetzt vorliegende Bemerkung eingearbeitet worden ist. Die meisten Nummern sind jetzt verständlich; die wenigen, die keine Bemerkungen zeigen, haben sich nicht aufklären lassen.

\indent Dem Abdruck des \textit{Tagebuchs} in den Werken ist eine photographische Nachbildung dieser wichtigen Urkunde vorangestellt worden, die alle Einzelheiten der Handschrift deutlich hervortreten läßt. Wir machen besonders auf die in den \textit{Vorbemerkungen} von \textsc{Klein} erwähnten Marken und auf die roten Unterstreichungen aufmerksam, die in der Nachbildung im Halbton wiedergegeben sind.
\par\noindent\hfill\textsc{Schlesinger}.
\ClosingRule
